\section{Strips}\label{sec:strip}

On two rows every rectangle is axis-aligned. Indeed, a tilted rectangle has no
horizontal side, so both sides at a lowest corner go strictly up, and the opposite
corner lies in a third row.

\begin{theorem}[strip theorem]\label{thm:strip}
Let $m\ge3$.
\begin{enumerate}[(a)]
\item A blackout of the $2\times m$ strip is valid if and only if at most one of its columns is full.
\item We have $f(2,m)=m+1$, and the maximum blackouts are those with one full column
and one black point in every other column.
\item There are $a_2(m)=m\cdot2^{m-1}$ maximum blackouts.
\end{enumerate}
For $m=2$ there is one rectangle, so $f(2,2)=4$ and $a_2(2)=1$.
\end{theorem}

\begin{proof}
(a) Two distinct rectangles share at most one column. If they
share one, it cancels, and the difference set is the four points of the other two
columns $p$ and $q$, the \emph{pair constraint} on $\{p,q\}$. If they share
no column, the difference set is the eight points of four columns, and it contains a
pair constraint. Every pair constraint occurs, since $m\ge3$ gives a third column $z$,
and the rectangles on $\{z,p\}$ and $\{z,q\}$ have exactly this difference set. By
Corollary~\ref{cor:hitting}, $S$ is valid if and only if $K$ meets every two columns,
that is, if and only if at most one column has no kept point.

(b) By (a) a valid $S$ has $|K|\ge m-1$, so $|S|\le m+1$. Equality holds exactly when
one column is full and every other column has one kept point, hence one black point.
Such blackouts are valid by (a).

(c) There are $m$ choices for the full column and two in each other column.
\end{proof}

So row $2$ of \seqnum{A399580} is \seqnum{A001787}$(m)=m2^{m-1}$ from $m=3$ on.

\begin{proposition}[orbits and kings]\label{prop:kings}
For $m\ge3$ the number $o_m$ of maximum blackouts of the $2\times m$ strip up to the row
reflection and the column reversal is $m\,2^{m-3}+[m\text{ odd}]\,2^{(m-3)/2}$. It equals
\seqnum{A322284}$(m-1)$, the number of ways to place $m-1$ nonattacking kings on a
$2\times(2m-2)$ board up to the symmetries of the rectangle. So $o_3,\ldots,o_8$ are
$4,8,22,48,116,256$.
\end{proposition}

\begin{proof}
By Burnside's lemma $4o_m$ is $a_2(m)$ plus the numbers of optima fixed by the reflection,
which is $0$ since it moves every singleton column, and by the reversal and the half turn. These fix none for even $m$, and
$2^{(m-1)/2}$ each for odd $m$, with a central full column.

Put $n=m-1$. An optimum is given by its full column $e$ and the rows $r_1,\ldots,r_n$ of its
other black points, from left to right. Two kings attack exactly when their columns are
equal or adjacent. So $n$ nonattacking kings on the $2\times2n$ board are given by columns
$c_1<\cdots<c_n$ with $c_{j+1}\ge c_j+2$ and any rows $r_1,\ldots,r_n$. The $n+1$ gaps
$c_1-1$, $c_{j+1}-c_j-2$ for $1\le j<n$, and $2n-c_n$ are nonnegative and sum to $1$. So
exactly one of them is $1$, say the $e$th, and $c_j=2j-1+[j\ge e]$. Hence the kings at
$(c_j,r_j)$ correspond bijectively to the optima. The reflection swaps the rows in both
pictures. The reversal sends $e$ to $n+2-e$ and reverses the rows, and on the board it
sends $c$ to $2n+1-c$, where $2n+1-c_{n+1-j}=2j-1+[j\ge n+2-e]$. So the bijection commutes
with both actions, and the orbits correspond.
\end{proof}
